
\documentclass[letterpaper, 11pt, sans]{moderncv}

\usepackage[utf8]{inputenc}
\usepackage[textwidth=7.5in, textheight=10in, footskip=20pt]{geometry}    % Width of the entire CV

\usepackage{moderntimeline}

\usepackage{color, graphicx}
\tlmaxdates{2006}{2025}              % Beggining and start of your timeline
\tlsetnotshadedfraction{0.2}

\usepackage{fancyhdr}
\usepackage{lastpage}
\rfoot{\footnotesize\color{gray}schloerke - \thepage  \hspace{1pt} / \pageref*{LastPage}}

\input{pre/style}

\input{pre/defs}

% Personal Information
\name{Barret}{Schloerke\vspace{-10pt}}
\title{{\large\emph{code[c(``architect'', ``manager'', ``philosopher'')]}}\newline{\large\emph{data[c(``wrangler'', ``explorer'', ``visualizer'')]}}}
% \address{Purdue University \\ Department of Statistics\\ 250 N. University St.\\}{West Lafayette, Indiana}{}
% \phone[mobile]{(123)~123~1234}
\email{schloerke@gmail.com}                % optional, remove / comment the line if not wanted
% \homepage{mdastro.wix.com/maria-luiza-dantas}    % optional, remove / comment the line if not wanted
\social[linkedin]{schloerke}  % optional, remove / comment the line if not wanted
% \social[twitter]{mdastro}  % optional, remove / comment the line if not wanted
\social[github]{schloerke}   % optional, remove / comment the line if not wanted
% \extrainfo{\emailsymbol \emaillink{mdantas.astro@gmail.com}}
% \quote{Dubitando ad veritatem parvenimus - \emph{Cicerone}}


\begin{document}


\makecvtitle

\vspace*{-8mm}

\section{Education}

\input{education/phd}

\input{education/masters}

\input{education/rice}

% \input{education/undergrad}
\tlcventry{2006/8}{2010/12}{Bachelor of Science in Computer Engineering}{Iowa State University}{Ames, IA}{3.77/4.0}{}

\section{Work Experience}


\tlcventry{2018/1}{2021/12}{\href{http://posit.co}{Posit (formerly RStudio)}}{}{Software Engineer}{}{
  \begin{itemize}
    \item Core developer of \href{https://shiny.posit.co/r/}{\texttt{Shiny}}, an R package that makes it easy to build interactive web apps straight from R.
    \item Architected, developed, and automated the nightly testing of 150+ Shiny applications using GitHub Actions over the Mac/Windows/Linux operating systems and 5x R versions.
    \item Successfuly converted core JavaScript code into TypeScript without minor bugs
    \item Helped develop and maintain over 45 R packages within the Shiny ecosystem
  \end{itemize}
}

\vspace{-10pt}

\input{work/metamarkets.tex}
% \vspace{-5pt}

% Healthy Birth, Growth, and Development Knowledge Integration
\tlrightcventry{2016/4}{2018/12}{Gates Foundation: \href{http://hbgdki.org}{HBGDki}}{Remote}{Tool Development Team}{}{
  \begin{itemize}
    \item Provided customized summary statistics in timely manner for 75+ datasets to feed into visualization applications
    \item Maintained data security while coordinating with five other data scientists
    \item Created and maintained shell R packages for standardized R package development
  \end{itemize}
}

% \vspace{-5pt}

\input{work/darpa}
% \vspace{-10pt}

\input{work/purdue_ta}

\input{work/hadley_ta}
  % \vspace{-10pt} % for lack of content
  % \vspace{-5pt}

\input{work/novartis}
\vspace{-5pt}


\newpage
\section{Honors and Awards}

\tlcventry{2012/8}{2017/4}{ \href{http://www.nsf.gov/grfp}{National Science Foundation Graduate Research Fellowship}}{}{}{}{
  \begin{itemize}
    \item Received a three-year fellowship (to be used within five years) to pursue graduate studies in a science based field
    \item Only awardee for Computational Statistics of the thirteen statistical recipients in 2012
    \item One of five recipients from Iowa State University in 2012
  \end{itemize}
}

\tlrightcventry{2015/2}{2015/3}{Visiting Scientific Researcher}{Monash University}{Melbourne, Victoria, Australia}{}{
  \begin{itemize}
    \item Presented latest research on \pkg{gqlr} to WOMBAT2016
    \item Collaborated with \href{http://numbat.space}{NUMBAT} on large data exploration tools
  \end{itemize}
}

% \tldatecventry{2015/2}{Visiting Scientific Researcher}{Fields Institute for
% Research in Mathematical Sciences}{Toronto, Ontario, Canada}{}{
%   \begin{itemize}
%     \item Spoke at `\href{http://www.fields.utoronto.ca/programs/scientific/14-15/bigdata/visualization/}{Visualization for Big Data: Strategies and Principles}' workshop
%   \end{itemize}
% }

% \bigskip

\tldatecventry{2011/6}{Best Dynamic / Interactive Visualization}{\href{http://datainsightsf.com/}{Data Insight (SF)}}{San Francisco, CA}{}{
\begin{itemize}
  \item Team of four worked together over the weekend to produce an interactive visualization
  \item Leaned on each other's skills to provide an application that highlighted our data and visualization talents
\end{itemize}
}

\tldatecventry{2011/4}{\href{http://www.nsf.gov/grfp}{National Science Foundation Graduate Research Fellowship}}{Honorable mention}{}{}{}

% \bigskip
%
% % \tldatecventry{2010/12}{Magna Cum Laude}{Graduated with GPA greater than 3.70}{}{}{}
%
% % Golden Keyboard &
% % Managed a database for Greek Week &
% % Spring 2010\\[0.5em]
%
% % Golden Key International Honour Society &
% % Being in the top 15\% of class &
% % Inducted in \newline
% % Fall 2008\\[0.5em]
%
% \tlcventry{2008/8}{2010/12}{Eta Kappa Nu Honor Society}{EE/CprE Honor Society for top 25\% of class}{}{}{}
%
% \tlcventry{2006/8}{2010/12}{Dean's List}{Achieving a 3.5 or above for a semester}{}{}{}
%
% \tlcventry{2006/8}{2010/12}{Engineering Merit Scholarship}{Achieving at least a 3.50 GPA in the College of Engineering}{}{}{}
%
% \tlcventry{2006/8}{2010/12}{President's Award for Competitive Excellence}{}{}{}{}






\section{Publications}

  \subsection{Journal Articles}

  \input{publication/lmmpar}

  \input{publication/boxland}

  \input{publication/generalized_pairs_plot}

  \tldatecventry{2012}{}{}{
    Hofmann H., D. Cook, C. Kielion, B. Schloerke, J. Hobbs, A. Loy, L. Mosley, D. Rockoff, Y. Huang, D. Wrolstad, and T. Yin ``\href{http://amstat.tandfonline.com/doi/abs/10.1198/jcgs.2011.3de}{Delayed, Canceled, on Time, Boarding… Flying in the USA.}'' \emph{Journal of Computational and Graphical Statistics} 20.2 (2012). Print.
  }{}{}

  \tldatecventry{2010}{}{}{
    Mosley, L., D. Cook, H. Hofmann, C. Kielion, and B. Schloerke. ``\href{http://chance.amstat.org/files/2010/12/Visually.pdf}{Monitoring the 2008 Election Visually.}'' \emph{Chance} 23.3 (2010). Print.
  }{}{}


    \subsection{Presentations}



    \tlrightdatecventry{2024/8}{}{}{
      Schloerke, B. \href{https://github.com/schloerke/presentation-2024-08-13-posit-shiny-data-frame}{``Editable data frames in Shiny for Python''} posit::conf. Posit. Aug. 2024. Speech.
      }{}{}

    \tlrightdatecventry{2024/6}{}{}{
      Schloerke, B. \href{https://www.youtube.com/watch?v=RcwvG7dtMqU}{``shinylive: Serverless Shiny apps''} YouTube. Posit. Jun. 2024. Video.
      }{}{}

    \tlrightdatecventry{2024/3}{}{}{
      Schloerke, B. \href{https://github.com/schloerke/presentation-2024-04-18-appsilon-shinylive}{``shinylive: Serverless Shiny apps''} Shiny Conference. Appsilon. May. 2024. Speech.
      }{}{}

    \tlrightdatecventry{2023/3}{}{}{
      Schloerke, B. \href{https://github.com/schloerke/presentation-2023-03-15-appsilon-nightly-testing}{``Lessons learned from testing 2500+ Shiny Apps every day''} Shiny Conference. Appsilon. May. 2023. Speech.
      }{}{}

    \tlrightdatecventry{2022/10}{}{}{
      Schloerke, B. \href{https://github.com/schloerke/presentation-2022-10-18-jnj22-shinytest2}{``shinytest2: Regression Testing for Shiny Applications''} Shiny Day. Johnson and Johnson. Oct. 2022. Speech.
      }{}{}

    \tlrightdatecventry{2022/7}{}{}{
      Schloerke, B. \href{https://github.com/schloerke/presentation-2022-05-26-rinpharma-shinytest2}{``shinytest2: Regression Testing for Shiny Applications''} rstudio::conf. RStudio. Forest Heights, MD. July. 2022. Speech.
      }{}{}


      \tlrightdatecventry{2022/6}{}{}{
      Schloerke \& Mostipak. \href{https://www.youtube.com/watch?v=BaQldCBZhCI}{``Programming Games with Shiny || Roll the Dice: with Quosures!''} YouTube. Posit. Jun. 2022. Video.
      }{}{}

      \tlrightdatecventry{2022/5}{}{}{
      Schloerke \& Mostipak. \href{https://www.youtube.com/watch?v=GNZY7gHSgzM}{``Programming Games with Shiny || Roll the Dice''} YouTube. Posit. May. 2022. Video.
      }{}{}

      \tlrightdatecventry{2022/5}{}{}{
      Schloerke \& Mostipak. \href{https://www.youtube.com/watch?v=87rpiuRyhaQ}{``Programming Games with Shiny || Guess the Number''} YouTube. Posit. May. 2022. Video.
      }{}{}

    \tlrightdatecventry{2022/5}{}{}{
      Schloerke \& Mostipak. \href{https://www.youtube.com/watch?v=sD39WAZo99A}{``Programming Games with Shiny || Dragon Realm''} YouTube. Posit. May. 2022. Video.
      }{}{}

    \tlrightdatecventry{2022/5}{}{}{
    Schloerke, B. \href{https://github.com/schloerke/presentation-2022-05-26-rinpharma-shinytest2}{``shinytest2: Regression Testing for Shiny Applications''} R in Pharma working group. May. 2022. Webinar.
    }{}{}

    \tlrightdatecventry{2022/4}{}{}{
      Schloerke, B. \href{https://github.com/schloerke/presentation-2022-04-27-appsilon-shinytest2}{``shinytest2: Regression Testing for Shiny Applications''} Shiny Conference. Appsilon. Apr. 2022. Speech.
      }{}{}

    \tlrightdatecventry{2022/2}{}{}{
      Schloerke, B. \href{https://www.youtube.com/watch?v=3gtk8uRrrL4}{``Maximize computing resources using future\_promise''} YouTube. Posit. Feb. 2022. Video.
      }{}{}

    \tlrightdatecventry{2021/8}{}{}{
      Schloerke, B. \href{https://www.youtube.com/watch?v=eHrzsIGY0so}{``Plumber: Asynchronous Route Execution with Barret Schloerke''} YouTube. Harvard Data Science Initiative. Aug. 2021. Webinar.
    }{}{}

    \tlrightdatecventry{2021/3}{}{}{
      Blaire \& Schloerke. \href{https://www.youtube.com/watch?v=J0Th2QRZ7Rk}{``Plumber: Asynchronous Route Execution with Barret Schloerke''} RStudio Webinar. Mar. 2021. Webinar.
    }{}{}

    \tlrightdatecventry{2021/1}{}{}{
      Schloerke, B. \href{https://github.com/schloerke/presentation-2021-01-rstudio-global-plumber-async}{``plumber + future: Async Web APIs''} rstudio::global. Jan. 2021. Speech.
    }{}{}

    \tlrightdatecventry{2020/8}{}{}{
      Schloerke, B. \href{https://shinydevseries.com/interview/ep014/}{``Episode 14: Shining a Light on learnr (Barret Schloerke Part 3)''} Shiny Developer Series with Eric Nantz. Aug. 2020. Podcast.
    }{}{}

    \tlrightdatecventry{2020/8}{}{}{
      Schloerke, B. \href{https://shinydevseries.com/interview/ep013/}{``Episode 13: Inside Plumber 1.0 (Barret Schloerke Part 2)''} Shiny Developer Series with Eric Nantz. Aug. 2020. Podcast.
    }{}{}

    \tlrightdatecventry{2020/8}{}{}{
      Schloerke, B. \href{https://shinydevseries.com/interview/ep012/}{``Episode 12: Barret Schloerke Part 1 (reactlog)''} Shiny Developer Series with Eric Nantz. Aug. 2020. Podcast.
    }{}{}

    \tlrightdatecventry{2020/04}{}{}{
      Schloerke, B. \href{https://github.com/schloerke/presentation-2019-11-11-shinymeta}{``Reactlog 2.0: Debugging the State of Shiny''} Remote Meetup: Shiny Deep Dive. Statistical Programming DC. Washington D.C. Apr. 2020. Speech.
    }{}{}

    \tlrightdatecventry{2019/11}{}{}{
      Schloerke, B. \href{https://github.com/schloerke/presentation-2019-11-11-shinymeta}{``shinymeta: Reproducible Shiny apps with shinymeta''} Autumn Biostatistics Conference US (ABACUS). GlaxoSmithKline. Phoenixville, PA. Nov. 2019. Speech.
    }{}{}

    \tlrightdatecventry{2019/5}{}{}{
      Schloerke, B. \href{https://github.com/schloerke/presentation-2019-05-02-shiny-reactlog-sparklyr}{``Debug Large Data Apps: shiny + reactlog + sparklyr''} Advanced R workshop. Northeastern University. Boston, MA. Mar. 2019. Speech.
    }{}{}

    \tlrightdatecventry{2019/1}{}{}{
      Schloerke, B. \href{https://github.com/schloerke/presentation-2019-01-18-reactlog}{``Reactlog 2.0: Debugging the State of Shiny''} rstudio::conf. RStudio. Austin, TX. Jan. 2019. Speech.
    }{}{}

    \tlrightdatecventry{2018/8}{}{}{
      Schloerke, B. \href{https://github.com/schloerke/presentation-2018-08-20-nasa-shiny-lightning/blob/master/-nasa-shiny-lightning.pdf}{``Shiny''} DATANAUTS. NASA. Washington, D.C. Aug. 2018. Speech.
    }{}{}

    \tlrightdatecventry{2017/4}{}{}{
      Schloerke, B. \href{https://github.com/schloerke/presentation-2017_04_14-cognostics}{``Cognostics: Metrics enabling detailed interactive visualization of big data''} CSESC2017. Purdue University, West Lafayette, IN. Apr. 2017. Speech.
    }{}{}

    \tlrightdatecventry{2017/1}{}{}{
      Schloerke, B. \href{https://github.com/schloerke/presentation-2017_01_26-tidyverse}{``tidyverse[c("magrittr", "dplyr", "tidyr")]''} Purdue Graduate Statistics Seminar. Math, West Lafayette, IN. Jan. 2017. Speech.
    }{}{}

    \tlrightdatecventry{2016/8}{}{}{
      Schloerke, B. \href{https://github.com/schloerke/presentation-2016_08_03_cognostics}{``Cognostics: Metrics enabling detailed interactive visualization of big data''} Joint Statistical Meeting. Chicago, IL. Aug. 2016. Speech.
    }{}{}

    \tlrightdatecventry{2016/6}{}{}{
      Schloerke, B. \href{https://github.com/schloerke/presentation-2016_06_30-ggduo}{``ggduo: Pairs plot for two group data''} UseR!2016. Stanford University, Stanford, CA. Jun. 2016. Speech.
    }{}{}

    \tlrightdatecventry{2016/6}{}{}{
      Schloerke, B. \href{https://slides.com/schloerke/tessera-src-2016-5-25}{``High Performance Computing for High-Resolution Analysis of Big Data and Small''} Spring Research Conference. Illinois Institute of Technology, Chicago, IL. Jun. 2016. Speech.
    }{}{}

    \tlrightdatecventry{2016/4}{}{}{
      Schloerke, B. \href{https://slides.com/schloerke/tessera-purdue-2016-4-7}{``Analysis and Visualization of Large Complex Data with Tessera''}. Graduate Statistics Seminar. Purdue, West Lafayette, IN. Mar. 2016. Speech.
    }{}{}

    \tlrightdatecventry{2016/3}{}{}{
      Schloerke, B. \href{http://slides.com/schloerke/tessera-isu-2016-3}{``Analysis and Visualization of Large Complex Data with Tessera''}. Graphics Research Group Seminar. Iowa State University, Ames, IA. Mar. 2016. Speech.
    }{}{}

    \tlrightdatecventry{2016/2}{}{}{
      Schloerke, B. \href{http://slides.com/schloerke/tessera-monash-2016}{``Analysis and Visualization of Large Complex Data with Tessera''}. NUMBAT Seminar. Monash University, Melbourne, VIC AUS. Feb. 2016. Speech.
    }{}{}

    \tlrightdatecventry{2016/2}{}{}{
      Schloerke, B. \href{https://github.com/schloerke/presentation-2016_02_18-graphql}{``GraphQL: An R server implementation''}. Wombat2016. Melbourne Zoo, Melbourne, VIC AUS. Feb. 2016. Speech.
    }{}{}

    \tlrightdatecventry{2015/10}{}{}{
      Schloerke, B. \href{https://github.com/schloerke/presentation-2015_10_20-web_scraping}{``Web Scraping with R''} American Credit Acceptance. Spartanburg, SC. Oct. 2015. Speech.
    }{}{}

    \tldatecventry{2015/2}{}{}{
      Schloerke, B. \href{https://github.com/schloerke/presentation-2015_02_24-trelliscope}{``Trelliscope: D\&R visualization tool''} Big Data Visualization. Fields Institute Toronto, CAN. Feb. 2015. Speech.
    }{}{}

    \tldatecventry{2014/10}{}{}{
      Schloerke, B. \href{https://github.com/schloerke/presentation-2014_10_21-ggplot2_spatial_statistics}{``ggplot2: displaying spatial and temporal data''} Spatial Statistics Seminar. Recitation, West Lafayette, IN. Oct. 2014. Speech.
    }{}{}

    \tldatecventry{2014/2}{}{}{
      Schloerke, B. \href{https://github.com/schloerke/presentation-2014_02_14-minimize-work-inefficiencies}{``Reducing Working Environment Inefficiencies''} Graduate Statistics Seminar. Haas, West Lafayette, IN. Feb. 2014. Speech.
    }{}{}

    \tldatecventry{2010/9}{}{}{
      Schloerke, B. ``helpr: Help for R.'' Working Group Meeting. Snedecor Hall, Ames, IA. Sept. 2010. Speech.
    }{}{}

    \tldatecventry{2010/8}{}{}{
      Schloerke, B. ``GGally: A Plot Matrix for All Variable Types.'' Joint Statistical Meeting 2010. Vancouver  Convention Center, Vancouver, BC Canada. Aug. 2010. Speech.
    }{}{}

    % \tldatecventry{2009/6}{}{}{
    %   Quart, G., B. Schloerke, and D. Gannon. ``Visualizing Large Data Sets: Housing Crisis.'' VIGRE. Rice  University, Houston, TX. June 2009. Speech.
    % }{}{}


\subsection{Lightning Talks}

  \tldatecventry{2014/6}{\href{https://github.com/schloerke/presentation-2014_06_27-Git}{``git''}}{}{Purdue Working Group}{}{}
  \tlcventry{2011/1}{2011/12}{``Sinartra: Simple Web Framework'', ``5 Tips for Running Node in Production'', ``Express: Node Router''}{}{San Francisco Bay Area}{}{}
  % \tlcventry{2008}{2011}{``Roxygen: Documentation for R'', ``helpr: Help for R'', and ``Electoral Towers''}{}{Iowa State University Working Group}{}{}



  \bigskip

% \newpage
\section{Projects}

  \subsection{Posit}
  \cvitem{Advisors:}{Joe Cheng and Winston Chang}

  \tlrightcventry{2023/12}{2024/1}{\href{https://github.com/posit-dev/py-shiny/tree/main/shiny/render/_data_frame.py}{\pkg{shiny.render.data\_frame}}}{\textnormal{Editable data frames within Shiny for Python applications}}{\emph{Posit}}{}{
  }
  \tlrightcventry{2023/12}{2024/1}{\href{https://github.com/posit-dev/py-shiny/tree/main/shiny/render/renderer/}{\pkg{shiny.render.renderer}}}{\textnormal{Extendable base class that creates renderers within Shiny for Python applications}}{\emph{Posit}}{}{
  }
  \tlrightcventry{2023/3}{0}{\href{https://github.com/posit-dev/py-shiny/tree/main/shiny/playwright}{\pkg{shiny.playwright}}}{\textnormal{Playwright testing controllers for every Shiny for Python UI element}}{\emph{Posit}}{}{
  }
  \tlrightcventry{2023/3}{0}{\href{https://github.com/posit-dev/py-shiny}{\pkg{shiny.pytest}}}{\textnormal{Pytest integration for testing Shiny for Python applications}}{\emph{Posit}}{}{
  }

  \tlrightcventry{2023/9}{0}{\href{https://github.com/posit-dev/r-shinylive}{\pkg{shinylive} for R}}{\textnormal{Export Shiny for R applications to a static folder to be hosted statically}}{\emph{Posit}}{}{
  }
  \tlrightcventry{2023/9}{2023/12}{\href{https://github.com/posit-dev/py-shinylive}{\pkg{shinylive} for Python}}{\textnormal{Export Shiny for Python applications to a static folder to be hosted statically}}{\emph{Posit}}{}{
  }
  \tlrightcventry{2023/9}{2024/5}{\href{https://github.com/posit-dev/shinylive}{\pkg{shinylive} assets}}{\textnormal{Assets required to run Shiny applications in the browser}}{\emph{Posit}}{}{
  }


  \tlrightcventry{2023/3}{0}{\href{https://github.com/posit-dev/py-htmltools}{\pkg{htmltools}}}{\textnormal{Bootswatch + Bootstrap 5 themes for Shiny for Python}}{\emph{Posit}}{}{
  }

  \tlrightcventry{2023/3}{0}{\href{https://github.com/posit-dev/py-shinyswatch}{\pkg{shinyswatch}}}{\textnormal{Bootswatch + Bootstrap 5 themes for Shiny for Python}}{\emph{Posit}}{}{
    }

  \tlrightcventry{2023/3}{0}{\href{https://github.com/posit-dev/py-shiny}{\pkg{shiny}}}{\textnormal{A web development framework for Python}}{\emph{Posit}}{}{
    }


  \bigskip

  \subsection{RStudio}
  \cvitem{Advisors:}{Joe Cheng and Winston Chang}

  \tlrightcventry{2019/12}{0}{\href{https://github.com/rstudio/shinycoreci}{\pkg{shinycoreci}}}{\textnormal{Regression testing for shiny-verse}}{\emph{RStudio}}{}{
  }
  \tlrightcventry{2018/4}{0}{\href{https://github.com/rstudio/shiny}{\pkg{shiny}}}{\textnormal{Easy interactive web applications with R}}{\emph{RStudio}}{}{
  }
  \tlrightcventry{2021/4}{2023/5}{\href{https://github.com/rstudio/htmltools}{\pkg{htmltools}}}{\textnormal{Tools for HTML}}{\emph{RStudio}}{}{
  }
  \tlrightcventry{2021/8}{2023/3}{\href{https://github.com/rstudio/shinytest2}{\pkg{shinytest2}}}{\textnormal{Testing for Shiny Appliations}}{\emph{RStudio}}{}{
  }
  \tlrightcventry{2021/7}{2022/9}{\href{https://github.com/rstudio/chromote}{\pkg{chromote}}}{\textnormal{Headless Chrome Web Browser Interface}}{\emph{RStudio}}{}{
  }
  \tlrightcventry{2018/11}{2022/1}{\href{https://github.com/rstudio/swagger}{\pkg{swagger}}}{\textnormal{Dynamically Generates Documentation from a 'Swagger' Compliant API}}{\emph{RStudio}}{}{
  }
  \tlrightcventry{2021/7}{2021/11}{\href{https://github.com/rstudio/webshot2}{\pkg{plumberDeploy}}}{\textnormal{Take screenshots of webpages from R}}{\emph{RStudio}}{}{
  }

  \tlrightcventry{2020/7}{2021/8}{\href{https://github.com/meztez/plumberDeploy}{\pkg{plumberDeploy}}}{\textnormal{Plumber Deployment}}{\emph{RStudio}}{}{
  }


  \tlrightcventry{2018/7}{2019/9}{\href{https://github.com/rstudio/gt}{\pkg{gt}}}{\textnormal{Easily Create Presentation-Ready Display Tables}}{\emph{RStudio}}{}{
    }

  \tlrightcventry{2018/6}{2022/12}{\href{https://github.com/rstudio/plumber}{\pkg{plumber}}}{\textnormal{Turn your R code into a web API.}}{\emph{RStudio}}{}{
  }

  \tlrightcventry{2018/6}{2020/10}{\href{https://github.com/rstudio/reactlog}{\pkg{reactlog}}}{\textnormal{Reactivity Visualizer for 'shiny'}}{\emph{RStudio}}{}{
  }

  \tlrightcventry{2018/1}{2018/8}{\href{https://github.com/rstudio/leaflet}{\pkg{leaflet}}}{\textnormal{R Interface to Leaflet Maps}}{\emph{RStudio}}{}{
  }


  \bigskip

  \subsection{Purdue University}
  \cvitem{Advisors:}{Dr. Ryan Hafen and Dr. William Cleveland}

  \tlrightcventry{2017/1}{2021/3}{\href{https://github.com/schloerke/autocogs}{\pkg{autocogs}}}{\textnormal{Automatic Cognostic Calculations}}{\emph{Purdue University}}{}{
  \begin{itemize}
    \item ``Automatically calculates cognostics for plot objects and list column plot objects. \pkg{autocogs} compliments \pkg{trelliscopejs}'s panel interactions by producing multiple cognositc values for the visualizations displayed''
    \item Generalized framework to produce consistent cognostics independent of visualization library utilized
  \end{itemize}
}


  \tlrightcventry{2015/11}{2017/6}{\href{https://github.com/schloerke/gqlr}{\pkg{gqlr}}}{\textnormal{\href{http://facebook.github.io/graphql/}{GraphQL} Server in R}}{\emph{Purdue University}}{}{
  \begin{itemize}
    \item ``R server implementation of \href{http://facebook.github.io/graphql/}{GraphQL}, a query language created by Facebook for describing complex data queries independent of the storage format.''
    \item \href{http://facebook.github.io/graphql/}{GraphQL} provides a complete and human readable description of the data in your data schema and gives clients the power to query only for exactly what they need.
    \item \pkg{gqlr} is a native R \href{http://facebook.github.io/graphql/}{GraphQL} implementation to be used with \href{https://facebook.github.io/relay/}{\texttt{Relay}} in \href{https://reactjs.org}{\texttt{React}}~javascript web applications
  \end{itemize}
}


  \tlrightcventry{2016/1}{2016/10}{\href{https://github.com/ggobi/ggally/}{\texttt{GGally::ggduo}}}{\textnormal{Generalized plot matrix for two-grouped data}}{\emph{Purdue University}}{}{
    \begin{itemize}
      \item Extension of ggpairs built on top of a ggmatrix object
      \item Pairs plot for two group data
      \item Readily useful for multiple time series analysis and canonical correlation analysis
      \item Foundation function for \href{http://ggobi.github.io/ggally/rd.html\#ggnostic}{`ggnostic'}, a model diagnostic plot matrix, and \href{http://ggobi.github.io/ggally/rd.html\#ggts}{`ggts'}, a time series plot matrix
      \item Funded by \href{https://github.com/rstats-gsoc/gsoc2016/wiki/ggduo:-pairs-plots-for-multiple-regression,-cca,-time-series}{Google Summer of Code} with \href{https://github.com/schloerke/gsoc-ggduo-final}{final report}
    \end{itemize}
  }

  \tlrightcventry{2015/1}{2016/1}{\href{https://github.com/ggobi/ggally/}{\texttt{GGally::ggmatrix}}}{\textnormal{Generalized plot matrix}}{\emph{Purdue University}}{}{
    \begin{itemize}
      \item Handles arbitrary plot columns and rows
      \item Used as display mechanism for all GGally plot matrix functions
    \end{itemize}
  }

  \tlcventry{2015/1}{2017/8}{\href{https://github.com/hafen/trelliscopejs}{\pkg{trelliscopejs}}}{\textnormal{Create Interactive Trelliscope Displays}}{\emph{Purdue University}}{}{
  \begin{itemize}
    \item ``Trelliscope is a scalable, flexible, interactive approach to visualizing data. This package provides methods that make it easy to create a Trelliscope display specification for \pkg{trelliscopejs}. High-level functions are provided for creating displays from within \pkg{dplyr}~or \pkg{ggplot2}~workflows. Low-level functions are also provided for creating new interfaces.''
    \item Ported \href{https://shiny.rstudio.com}{Shiny}-based \pkg{trelliscope}~R package to be built with \href{http://reactjs.com}{React} framework to increase interaction speed
    \item Integrated with \pkg{ggplot2}~objects to seamlessly produce \pkg{trelliscopejs} applications
  \end{itemize}
}


  \input{pkg/rbokeh}

  \tlcventry{2013/4}{2013/8}{\href{https://github.com/schloerke/Bitcoin-Transaction-Network-Extraction}{Bitcoin}}{DARPA XDATA}{}{}{
    \begin{itemize}
      \item Brought previous transaction exporter up to protocol
      \item Formatted escrow transactions to be matched easier
      \item Made ``None'' keys uniquely identifiable, rather than being treated the same
      \item Matched some ``None'' keys to valid output keys
      \item Updated User exporter to include all information available
    \end{itemize}
  }

\bigskip

  \subsection{Metamarkets}

  \tlcventry{2011/5}{2012/7}{\href{http://github.com/vogievetsky/DVL}{DVL}}{\textnormal{Data Visualization Legos}}{\emph{Metamarkets}}{}{
  \begin{itemize}
    \item Functionally reactive JavaScript library to bind data to functions and visualizations
  \end{itemize}}

  \tlcventry{2012/1}{2012/4}{\href{http://github.com/schloerke/yalog}{yalog}}{\textnormal{Yet Another Logger}}{\emph{Metamarkets}}{}{
  \begin{itemize}
    \item Request logger to allow for personalized logging messages within different log levels
  \end{itemize}
  }



  \bigskip

\subsection{Rice University}
  \cvitem{Advisor:}{Dr. Hadley Wickham}

  \tlcventry{2010/4}{2010/8}{\href{http://github.com/hadley/helpr/}{helpr}}{Help for R}{}{}{
    \begin{itemize}
      \item Betters Friendly HTML documentation with links to other packages, function aliases, and function sources.  Finding information may be done with the comprehensive search bar.
    \end{itemize}
  }
  % Helpr is an R package that betters friendly HTML documentation. With links to other packages, function aliases, and function sources, finding information is a click away. Using the comprehensive search bar, searching across all R packages is quick and effortless.
  %
  % The heart of helpr is hosted locally. No internet is required to display all of the documentation, while functionality of the search bar, RSS feed, and comment system requires an internet connection.
  %
  % The development of helpr was made possible with generous support from Revolution Analytics and direction from Hadley Wickham.


  % using 2010 dates for pretty alignment
  \tlcventry{2010/4}{2010/8}{VIGRE Program}{}{}{}{
    \begin{itemize}
      \item Performed exploratory analysis on the housing crisis
      \item Found, cleaned, and explored multiple data sets
      \item Created all outputs reproducible from scratch
    \end{itemize}
  }
  % VIGRE is a program sponsored by the National Science Foundation to carry out innovative educational programs in which research and education are integrated and in which undergraduates, graduate students, postdoctoral fellows, and faculty are mutually supportive.  During my stay at Rice University, I worked daily with two other undergraduate students, one graduate student, and our professor, Dr. Hadley Wickham.
  %
  % Our group project was to perform exploratory analysis on the housing crisis.  I found, cleaned, and explored multiple data sets.  The R programming language was used to clean the data.  To stay up-to-date with the other group members, Git, a repository language that evolved from SVN, was used constantly throughout the program.  One big requirement, which was followed during the project, was reproducibility.  Anyone could start  by retrieving our code and reproduce everything that we have already produced by sourcing the files we have made.
  %
  % Near the end of the internship, my undergraduate group gave a 20-minute speech on our findings, what we were currently working on, and where the group was heading in the future.



  \bigskip

\subsection{Iowa State University Statistics Research}
  \cvitem{Advisors:}{Dr. Dianne Cook and Dr. Heike Hofmann}
  % Starting in the Summer of 2007, I worked under Dr. Dianne Cook and Dr. Heike Hofmann, ISU Statistics professors. Since then, I have explored different topics in graphics. Each project has touched on different areas. (Starting from most recent).

  \input{pkg/ggally}

  \tlcventry{2010/1}{2010/12}{\href{http://github.com/ggobi/cranvas/}{cranvas}}{New Plotting Canvas for R}{}{}{
    \begin{itemize}
      \item New plotting canvas for R.  Produces more than a million points in less than a second.
      \item Dr. Wickham, Dr. Lawrence, Dr. Hoffman, Dr. Cook, Yihui Xie, Tengfei Yin, Marie Vendettuoli, Barret Schloerke
    \end{itemize}
  }


  \tlcventry{2009/8}{2009/12}{\href{http://github.com/ggobi/DescribeDisplay}{DescribeDisplay}}{Reproduce GGobi images in R}{}{}{
    \begin{itemize}
      \item Produce publication quality graphics from GGobi's describe display plugin.
    \end{itemize}
  }
  % Produce publication quality graphics from output of GGobi's describe display plugin

  \tlcventry{2009/8}{2009/12}{\href{http://github.com/ggobi/tourr}{tourr}}{Geodesic interpolation and basis generation functions}{}{}{
    \begin{itemize}
      \item Implements geodesic interpolation and basis generation functions that allow you to create new tour methods in R.
      \item Dr. Wickham, Dr. Cook, Barret Schloerke
      \item Iowa State University and Rice University
    \end{itemize}
  }
  % Implements geodesic interpolation and basis generation functions that allow you to create new tour methods from R.


  \tlleftcventry{2009/1}{2009/5}{\href{http://stat-computing.org/dataexpo/2009/posters/hofmann-cook.pdf}{American Statistical Association Data Expo}}{Delayed, Cancelled, On-Time, Boarding ... Flying in the USA}{}{}{
    \begin{itemize}
      \item Displayed airport delays caused by airport layout using multiple applications.
      \item Using Google's KML language, R and wget, we produced Google Earth images with the runways highlighted.
      \item Iowa State University working group finished 2nd in the \href{http://stat-computing.org/dataexpo/2009/posters/}{world competition}.
    \end{itemize}
  }
  % \item{2009 ASA Data Expo:} Delayed, Cancelled, On-Time, Boarding ... Flying in the USA
  %
  % The 2009 ASA Data Expo aim was to produce a graphical summary of important features.  Areas that I worked on dealt with which carriers were the best and why they were the best, which carriers performed the least amount of ``ghost flights'' (empty flights) and whether or not the layout of the airport was correlated with delays.
  %
  % Displaying the airport delays caused by airport layout used multiple applications.  Using Google's KML language, R and wget, Chris and I were able to produce Google Earth Images with the runways highlighted.  Following this, we produced images of the wind tendencies at each airport and matched them up with the corresponding airport runway direction.
  %
  % The group finished 2nd in the \href{http://stat-computing.org/dataexpo/2009/posters/}{world competition}.



  \tlleftcventry{2008/8}{2008/12}{\href{http://isu_research.schloerke.com/Election/}{Election Data}}{Explore new ways to show the Electoral College for the 2008 Presidential Election}{}{}{
    \begin{itemize}
      \item Created a custom cartogram of the United States with area proportional to the electoral vote of each state.
      \item Produced the `Electoral Vote Towers' (inspired by New York Times 2000 Presidential Race) that quickly portrayed the electoral vote status.
    \end{itemize}
  }

  % This project focused on the electoral votes of each state. A cartogram of the United States with the area proportional to the electoral vote of each state and the `Electoral Vote Towers' were the two main graphics that I helped develop.
  %
  % The cartogram shows voting pattern in a way that is more representative of the electoral votes while maintaining some geographic relations. The cartogram is block-constructed, coded, and imported into R. It demonstrates the states' area in proportion to its electoral votes. For example, Montana only has an area of three votes, while Rhode Island has an area of four votes even though it is much smaller geographically.  (Votes are determined by the number of representatives within the House and Senate, not by geographic area).
  %
  % The `Electoral Vote Towers' was inspired by a New York Times display during the 2000 presidential election titled ?Building an Election Victory?. These displays show the sum of the electoral votes for each candidate. The shape comes from each state's percent difference (horizontal) and electoral votes (vertical). While the colors may be redundant, it still helps categorize the states' voting percentages. The goal of the candidates is to reach the 270 votes needed to win the presidency which corresponds to the towers building up to the 270th line. Each of the states that are still a toss-up have their relevant statistics reported on the candidate's side to which they are currently leaning.




  \tlleftcventry{2008/8}{2008/12}{\href{http://isu_research.schloerke.com/DNA/index.html}{DNA}}{Viewing DNA sequence data}{}{}{
    \begin{itemize}
      \item Extension of the polytope ranking.
      \item Created combinations of DNA sequences with a length of four which were displayed in a grand tour.
    \end{itemize}
  }

% Dr. Michael Lawrence and Fred Hutchinson Cancer Research Center were working on a DNA data set and needed a way to view the distribution of DNA sequence samples relative to the possible combinations. The output needed to have a feasible way to tell trends while still showing as much data as possible.
%
% The original goal was to make outputs for a sequence of 6 digits but was later reduced to 4 digits when some constraints were imposed. The project was an extension of the polytope ranking from Spring 2008. It developed proportionally sized spheres at each vertice/combination with points jittered on the surface of the sphere. This makes the most common sequences `pop' out at the viewer.

  \tlleftcventry{2008/1}{2008/5}{Polytope Ranking}{Visualizing ranked data}{}{}{
    \begin{itemize}
      \item Created a model that would connect a ranked data set.
      \item Added jittered data points to each combination with the distance relative to the amount of points.
    \end{itemize}
  }

  \tlleftcventry{2007/8}{2007/12}{\href{http://isu_research.schloerke.com/FaceOff/}{FaceOff}}{`Lines and Dots' in the fourth dimension}{}{}{
    \begin{itemize}
      \item Game created to calibrate users into viewing and thinking in the fourth dimension.
      \item Game is a combination of  ``Connect 4'', ``Tic Tac Toe'', and ``Dots and Lines.''
    \end{itemize}
  }

% FaceOff is a game that is designed to help calibrate your eyes into viewing and mind into thinking in the fourth dimension. In 2004, the game arose from discussions between David Bulger and Dr. Dianne Cook as a first attempt to grow a 2-D game into higher dimensions. It has been further developed by Spencer Bradley, Dr. Hadley Wickham , Dr. Heike Hofmann and myself.
%
% The goal of the game is to connect a square in a 4-D cube with all four corners painted with your color. The game is a combination of ``Connect 4'', ``Tic Tac Toe'', and ``Dots and Lines'', only it is played in 4-D.  It is similar to `Connect 4' in that you need all four dots of your color, `Dots and Lines' in that you make a square, and the strategies from `Tic Tac Toe'.

% The FaceOff \href{http://www.public.iastate.edu/~bigbear/FaceOff/index.html}{website} has a video and PDF tutorial.


\tlleftcventry{2007/8}{2007/11}{IEEE Infovis Competition}{2007 Infovis visualization competition}{}{}{
  \begin{itemize}
    \item Internet Movie Data Base (IMDB) competition.
    \item Worked on the genres of the database using the ``Many Eyes'' web application.
    \item Helped edit and produce the movie that the working group submitted.
  \end{itemize}
}

% The 2007 Infovis competition dealt with Internet Movie Data Base (IMDB). I contributed by doing work on the genres of IMDB's movie database through the use of Many Eyes.  The link for the movie that I helped edit is located on Dr. Hadley Wickham's \href{http://had.co.nz/infovis-2007/}{website}.

% \subsection{Honors Research}
% \begin{itemize}
%   \item Modernize a strain gauge and monitor into a LabView program.
%   \item Strain gauge was used to determine friction on the material.
%   \item Supervised by Dr. Sriram Sundararajan.
% \end{itemize}

% The Freshman Honors program provides a wonderful opportunity for freshman undergrad students to participate in research.  I participated under Dr. Sriram Sundararajan, a Mechanical Engineer professor,  on micro/nanoscale tribology.  The goal of the project was to take the old strain gauge and strain gauge monitor and turn it into a modern-day strain gauge system.  Another program member and I learned how to use LabVIEW and then used our knowledge to program a system to interpret the readings from the strain gauge attached to the system.  The program that we made allowed the user to record the strain caused by the friction from the spinning material over time.  The program also calculated the coefficient of friction at each time.  We ran several test runs on the new system to check it.

\tlleftcventry{2007/1}{2007/12}{\href{http://schloerke.github.io/geozoo/}{GeoZoo}}{A library of high dimension geometric objects}{}{}{
  \begin{itemize}
    \item Functions to produce cubes, spheres, simplices, polyhedra, polytops, tori, and mobius-like objects.
    \item Contains a new addition to the high dimension polytopes, a Hyper Ring Torus
  \end{itemize}
}

%  \href{http://streaming.stat.iastate.edu/~dicook/geometric-data/}{website} was created to help the viewer's eyes calibrate to high dimensions. GeoZoo contains many data sets of high dimensional geometric shapes. This includes cubes, spheres, simplices, polyhedra, polytopes, tori, and mobius like objects. GeoZoo contains a new addition to the high dimension polytopes, a Hyper Ring Torus. This work was jointly supervised by Dr. Hadley Wickham and Dr. Dianne Cook.
%
% From my work on GeoZoo, I submitted a paper on the objects to The R Journal.



% \newpage

% \newpage

  \section{Service}

  \input{service/dubh}
  \input{service/r4ds}
  \input{service/gso}

  \tlcventry{2014/8}{2016/4}{StatCom Assistant}{}{}{}{
    \begin{itemize}
      \item Coordinated and put on multiple events
      \item Helped in production of StatCom banner
    \end{itemize}
  }

  \input{service/stat_mentor}

\section{Technical}

\vspace{-5pt}
\cventry{Languages}{}{}{}{}{
\vspace{-2em}
  \begin{itemize}
    \item Expert: R, Python, JavaScript, TypeScript, Node.js, GraphQL, Markdown, HTML, \LaTeX, JSON, YAML
    \item Moderate: MySQL, JSX, Regular Expressions, Bash, CSS, C
    % \item Beginner: C++, Java, VBA
  \end{itemize}
}

\input{extra/systems}


% \input{extra/coursework}

% \newpage
% \section{References}
% \input{extra/references}

\end{document}
